% Chapter 1: Introduction
\chapter{Introduction}
\label{ch:introduction}

\section{Background and Motivation}
\label{sec:background}

The logistics sector is one of Hong Kong's most economically significant industries, second only to financial services \autocite{chan2025hkindustrial}. As one of the world's most densely urbanised cities and a major logistics hub in the Asia--Pacific region, delivery operations must contend with congested road networks, compact geographic areas, and demanding service-level expectations. Small and medium-sized enterprises (SMEs), which form the backbone of the local logistics market, are often the most affected; unlike large corporations with access to proprietary fleet-management platforms, SMEs commonly rely on fragmented tools for dispatching, navigation, and proof of delivery \autocite{bhattacharjee2024digitizing}.

This fragmentation manifests in several concrete problems. Dispatchers lack a single unified interface for assigning tasks, monitoring driver progress, and reviewing delivery evidence. Drivers, in turn, operate across multiple disjointed applications: tasks arrive via one channel, navigation occurs in another, and proof-of-delivery (PoD) reporting is handled separately. The result is suboptimal routing, communication delays, and increased operating costs. Beyond commercial implications, poorly coordinated operations can compromise worker welfare: drivers navigating unfamiliar locations under tight schedules and without adequate route guidance face elevated stress and safety risks \autocite{apdullah2021occupational}.

Meanwhile, the rapid maturity of large language models (LLMs) has opened new avenues for operational decision support. Recent research demonstrates that generative AI can meaningfully reduce cognitive load and accelerate information retrieval in enterprise contexts \autocite{brynjolfsson2023generative}, yet its application within last-mile logistics remains largely unexplored. The convergence of these factors (an industry need for unified digital tools, the limitations of existing commercial solutions, and the emergence of practical AI capabilities) motivated the development of OrderVista.

\section{Problem Statement}
\label{sec:problem}

Existing delivery management tools either target large-scale enterprises with complex enterprise resource planning (ERP) integrations, or provide only basic route-planning features without task management, driver coordination, or AI-powered assistance. No publicly available platform combines unified task lifecycle management, intelligent route sequencing, multi-tenant security, proof-of-delivery workflows, and an AI assistant within a single, affordable system tailored for logistics SMEs. This project addresses that gap.

\section{Aim}
\label{sec:aim}

The aim of this project is to design, implement, and evaluate \textbf{OrderVista}, an AI-enhanced logistics dispatch platform that unifies administrative dispatching and driver-side field operations within a single cross-platform application, with the goal of reducing delivery distance, improving operational coordination, and demonstrating the practical utility of LLM-based assistance in last-mile logistics.

\section{Objectives}
\label{sec:objectives}

The project defines fourteen objectives, grouped into core platform objectives and advanced optimisation objectives.

\subsection{Core Platform Objectives}
\begin{enumerate}
  \item \textbf{Unified platform development:} Deliver a cohesive system serving both administrative and driver users, deployable across web and mobile from a single Flutter codebase.
  \item \textbf{Task management:} Implement end-to-end workflows for task creation, assignment, tracking, and completion, including CSV bulk import with validation.
  \item \textbf{Route sequencing:} Develop a baseline nearest-neighbour sequencing algorithm to order delivery tasks by proximity, with a map-based visual preview.
  \item \textbf{Proof of delivery:} Enable drivers to submit photographic evidence and notes upon task completion, with administrative review capabilities.
  \item \textbf{AI assistant integration:} Provide natural-language features for task summarisation and conversational queries about delivery operations, powered by the DeepSeek R1 model \autocite{deepseek2025r1}.
  \item \textbf{Multi-tenant security:} Enforce data isolation between logistics companies using PostgreSQL Row-Level Security (RLS) policies \autocite{postgresql2026rls}.
\end{enumerate}

\subsection{Advanced Optimisation Objectives}
\begin{enumerate}[resume]
  \item \textbf{Driver recommendation:} Recommend task-to-driver assignments based on availability, GPS proximity, and a composite performance score incorporating historical delivery metrics.
  \item \textbf{Customer communication:} Support customer-facing status queries and automated delivery notifications.
  \item \textbf{Agentic AI:} Extend the AI assistant to trigger structured workflow actions (e.g., sending notifications, reassigning tasks) under explicit human-oversight safeguards.
  \item \textbf{Real-time communication:} Enable live chat between administrators and drivers for operational coordination.
  \item \textbf{Dynamic priority recalculation:} Automatically adjust task priority based on time-window urgency and accumulated delays.
  \item \textbf{Mid-shift re-sequencing:} Re-optimise the task sequence from the driver's current GPS position after each task completion.
  \item \textbf{Fleet-wide distance minimisation:} Assign drivers to task clusters to minimise total fleet driving distance rather than optimising individual routes in isolation.
  \item \textbf{Performance-aware driver scoring:} Maintain per-driver performance metrics (on-time rate, tasks per hour, completion rate) and incorporate them into assignment recommendations.
\end{enumerate}

All fourteen objectives were achieved during the project. Chapters~\ref{ch:implementation} and~\ref{ch:testing} provide implementation details and evaluation evidence, respectively.

\section{Scope and Limitations}
\label{sec:scope}

\paragraph{In scope.}
The project encompasses a web-based administrative portal, an Android mobile driver application, a cloud-hosted backend (Supabase: PostgreSQL, Auth, Storage), external API integrations (Mapbox for geocoding and static maps, OpenRouter/DeepSeek for AI), multi-tenant data architecture with RLS, baseline and advanced route sequencing algorithms, proof-of-delivery capture and review, AI-powered summarisation and agentic actions, and a real-time admin--driver chat system.

\paragraph{Out of scope.}
The following items fall outside the project boundary: continuous real-time GPS vehicle tracking (telemetry is event-based), advanced VRP solvers such as genetic algorithms or branch-and-bound methods, Google Play Store publication, deep integrations with third-party ERP or e-commerce platforms, and large-scale field trials with commercial logistics companies (testing uses representative synthetic data within the university network environment).

\paragraph{Constraints.}
Development was conducted by a single developer within an academic timeline. Testing was primarily performed using emulators, web browsers, and limited physical-device trials. External API usage is subject to provider rate limits and availability on free-tier plans. No proprietary logistics datasets were available; evaluation therefore uses manually curated test data.

\section{Report Structure}
\label{sec:report-structure}

The remainder of this report is organised as follows:
\begin{itemize}
  \item \textbf{Chapter~\ref{ch:literature}} reviews relevant literature on vehicle routing, multi-tenant security, offline-first architectures, and AI in logistics.
  \item \textbf{Chapter~\ref{ch:design}} presents the system design, including architecture, data model, security approach, and critical workflows.
  \item \textbf{Chapter~\ref{ch:implementation}} details the methodology, implementation of each subsystem, and key algorithms.
  \item \textbf{Chapter~\ref{ch:testing}} reports testing procedures, performance measurements, and evaluation results.
  \item \textbf{Chapter~\ref{ch:discussion}} provides a critical discussion of achievements, alternatives considered, and limitations.
  \item \textbf{Chapter~\ref{ch:conclusion}} concludes the report and outlines directions for future work.
\end{itemize}
